% 中文摘要
\begin{cnabstract}

随着人工智能技术的飞速发展，图像识别作为计算机视觉领域的核心任务之一，在医疗诊断、自动驾驶、安防监控等领域发挥着越来越重要的作用。传统的图像识别方法依赖于人工设计的特征提取器，难以应对复杂多变的实际场景。近年来，深度学习技术的兴起为图像识别带来了新的突破，其中卷积神经网络（CNN）已成为该领域的主流方法。

本文针对图像识别中的关键技术问题展开研究，主要工作包括以下几个方面：首先，系统分析了卷积神经网络的基本原理和网络结构设计，探讨了残差网络（ResNet）、密集连接网络（DenseNet）等经典架构的优缺点；其次，针对现有模型参数量大、计算复杂度高的问题，提出了一种轻量级的改进网络结构，在保证识别精度的前提下显著降低了模型参数量和计算成本；最后，在CIFAR-10和ImageNet两个公开数据集上进行了大量实验验证，实验结果表明本文提出的方法相比基线模型具有更好的性能表现。

本文的研究工作为轻量化图像识别模型的设计提供了新的思路，对于推动深度学习在移动端和边缘设备上的应用具有重要的理论意义和实践价值。

\cnkeywords{图像识别；深度学习；卷积神经网络；残差网络；轻量级网络}
\end{cnabstract}

% 英文摘要
\begin{enabstract}

With the rapid development of artificial intelligence technology, image recognition, as one of the core tasks in the field of computer vision, plays an increasingly important role in medical diagnosis, autonomous driving, security surveillance and other fields. Traditional image recognition methods rely on manually designed feature extractors, which are difficult to cope with complex and changing real-world scenarios. In recent years, the rise of deep learning technology has brought new breakthroughs to image recognition, among which Convolutional Neural Networks (CNN) have become the mainstream method in this field.

This paper conducts research on key technical issues in image recognition. The main work includes the following aspects: Firstly, this paper systematically analyzes the basic principles of convolutional neural networks and network structure design, exploring the advantages and disadvantages of classic architectures such as Residual Network (ResNet) and Dense Network (DenseNet). Secondly, to address the problems of large parameter quantity and high computational complexity of existing models, a lightweight improved network structure is proposed, which significantly reduces model parameters and computational cost while ensuring recognition accuracy. Finally, extensive experimental verification is conducted on CIFAR-10 and ImageNet public datasets. Experimental results show that the proposed method has better performance compared with baseline models.

The research work in this paper provides new ideas for the design of lightweight image recognition models, and has important theoretical significance and practical value for promoting the application of deep learning on mobile terminals and edge devices.

\enkeywords{Image Recognition; Deep Learning; Convolutional Neural Network; Residual Network; Lightweight Network}
\end{enabstract}
